\begin{solapartat}
La força que fa un camp magnètic sobre una partícula carregada en moviment és $\vec{F} = q\,\vec{v}\wedge\vec{B}$ \punts{0,2}\unskip. Com que ara la velocitat i el camp magnètic són perpendiculars, el mòdul de la força és simplement $F = q\,v\,B$. \punts{0,1} L'electró i el protó tenen la mateixa càrrega en valor absolut de forma que el mòdul de la força que fa el camp és el mateix
\[ F = (1,60\cdot 10^{-19})(5)(2\times 10^{-3}) \equiv 1,60\cdot 10^{-21}\ \mathrm{N}\,. \quad\punts{0,2} \]
Per altra banda, el neutró no té càrrega elèctrica i per tant la força que actua sobre ell és $F = 0\ \mathrm{N}$. \punts{0,1}

Atès que la càrrega del protó és positiva, la força que fa el camp es dirigeix sobre l'eix X i en sentit positiu. \punts{0,1} El protó descriu llavors un moviment circular en sentit horari vist des del semieix positiu Z. \punts{0,1} Per altra banda, la càrrega de l'electró és negativa i el moviment que realitza és també circular, però en sentit contrari al que realitza el protó. \punts{0,1} El neutró segueix una trajectòria rectilínia a velocitat constant segons el sentit positiu de l'eix Y. \punts{0,1}
\end{solapartat}

\begin{solapartat}
El camp magnètic creat per un fil infinit de corrent sobre un punt que es troba a una distància $r$ de l'eix del fil és
\[ B = \frac{\mu_0 I}{2\pi r}\,. \]
Per tal que el protó descrigui un moviment rectilini, cal que el camp magnètic total que actua sobre ell sigui 0. En el present cas, això vol dir que el fil ha de fer un camp $\vec{B}_{fil} = -2,00\cdot 10^{-3}\,T\,\hat{k}$. \punts{0,4}

Pel que fa al valor del corrent, el trobem imposant la condició
\[ \frac{\mu_0 I}{2\pi r} = 2,00\cdot 10^{-3}\,, \quad\punts{0,1} \]
o el que és el mateix
\[ I = \frac{4\pi\, 10^{-3}\, 3\cdot 10^{-3}}{4\pi\cdot 10^{-7}} = 30,0\ \mathrm{A}\,. \quad\punts{0,5} \]
\end{solapartat}
